% ============================================================
%  Chapter 1: Introduction
%  File: chapters/01_introduction.tex
%  Bibliography: chapters/refs_01_intro.bib
% ============================================================

\chapter{Introduction}
\label{ch:introduction}

\section{Motivation}

The design of artificial neural networks has historically been guided by
human intuition. From the perceptron model of \citet{rosenblatt1958perceptron}
to modern deep architectures \citep{lecun2015deep}, researchers have manually
specified the number of layers, the connectivity patterns, the activation
functions, and the learning rules. This top-down approach has produced
remarkable results, but it has also imposed structural constraints that may
not be optimal for every task.

A central example of such a constraint is the layered architecture itself.
Since the backpropagation algorithm was popularized by
\citet{rumelhart1986learning}, nearly all successful neural networks have
been organized into discrete layers. This design is computationally
convenient---it allows efficient matrix multiplication and gradient
computation---but it also imposes a rigid separation between
``input,'' ``hidden,'' and ``output'' units. In a biological brain, by
contrast, there are no layers. Neurons form a dense, recurrent graph in
which any neuron may connect to any other, and structure emerges through
developmental and evolutionary processes rather than explicit design.

This discrepancy raises a fundamental question:

\begin{quote}
\emph{Can a neural network discover its own structure from first
principles, without a designer specifying layers, connectivity, or even
the number of nodes?}
\end{quote}

The present work answers this question in the affirmative, through a
minimal architecture we call \modnet{}.

\section{Background}

\subsection{The XOR Problem and the Limits of Linear Models}

The limitations of fixed, single-layer architectures were formalized by
\citet{minsky1969perceptrons}, who proved that a single-layer perceptron
cannot compute the XOR function, nor the parity function for any number of
inputs. This result, often referred to as the ``XOR problem,'' is not merely
a technical curiosity. It reflects a deeper mathematical fact: linear
threshold units can only separate linearly separable classes, and parity is
not linearly separable for any $n \geq 2$.

The XOR problem had a profound impact on the field. It contributed to the
first ``AI winter'' and demonstrated that architectural choices---in
particular, the presence or absence of hidden layers---are not merely
implementation details but determine the computational power of the model.
It was not until the rediscovery of backpropagation by
\citet{rumelhart1986learning} that multi-layer networks became trainable,
and the XOR problem was finally solved by adding a single hidden layer.

\subsection{Evolutionary Approaches to Network Design}

An alternative to gradient-based training is to evolve the network topology
itself. The NeuroEvolution of Augmenting Topologies (NEAT) algorithm,
introduced by \citet{stanley2002evolving}, was a landmark contribution in
this direction. NEAT starts with a minimal network and incrementally adds
nodes and connections, using a principled crossover mechanism and
speciation to protect structural innovation. It demonstrated that evolving
topology alongside weights can significantly accelerate learning on
reinforcement learning benchmarks.

NEAT was later extended by HyperNEAT \citep{stanley2009hypercube}, which
uses an indirect encoding called a Compositional Pattern-Producing Network
(CPPN) to generate connectivity patterns with geometric regularity.
HyperNEAT enabled the evolution of much larger networks than direct
encoding, but it still relies on a predefined substrate geometry, typically
a grid of neurons.

A parallel line of work explored growing neural networks with incremental
learning. The Cascade-Correlation architecture of
\citet{fahlman1990cascade} begins with a minimal network and adds hidden
units one by one. The Growing Neural Gas algorithm of
\citet{fritzke1995growing} uses competitive Hebbian learning to build
topological representations incrementally. Both approaches demonstrate the
power of growth, but neither fully addresses the problem of recurrence,
nor the possibility of a fully layer-free topology.

\subsection{Modularity in Neural Networks}

Modularity---the decomposition of a system into cohesive, loosely coupled
components---is a well-established principle in software engineering
\citep{parnas1972modularity}. \citet{alho2024modularity} conducted a
systematic review of $86$ studies on modular neural networks and found that
nearly two-thirds report improvements in task accuracy compared to
monolithic solutions.

In the context of neuroevolution, \citet{khare2005coevolutionary} proposed
a co-evolutionary model for automatic problem decomposition, in which a
population of modules is evolved alongside a population of complete
systems that combine modules. They demonstrated that if a particular task
decomposition is beneficial, it can be evolved through this mechanism.
However, their model requires the designer to specify the number of modules
and the interface between them.

\citet{landassuri2012evolution} investigated the emergence of modularity
in evolved neural networks and found that pure-modular architectures emerge
at low connectivity values. These results suggest that modularity is not
merely a design choice but can be an emergent property of evolutionary
search.

In the present work, we deliberately \emph{abandon} explicit modular
structure. Our architecture has no modules at all. This is a stronger
demonstration of emergence: the properties we observe---recurrence,
sparsity, cross-scale integration---must arise entirely from the
evolutionary process, without any modular scaffold to support them.

\subsection{Sparse Connectivity}

Biological neural networks are highly sparse: the probability that any two
neurons are connected is on the order of $10^{-6}$ \citep{lecun2015deep}.
This sparsity is not accidental; it is a consequence of metabolic
constraints and is believed to be essential for efficient information
processing \citep{chklovskii2004wiring}.

In artificial neural networks, sparsity has been studied primarily as a
compression technique. The Sparse Evolutionary Training (SET) method of
\citet{mocanu2018scalable} evolves an initial sparse topology into a
scale-free topology during training, allowing the training of networks
with over one million neurons on commodity hardware. However, these methods
typically begin with a predefined network architecture and prune or
sparsify it. They do not address the question of how sparsity and
connectivity structure emerge from first principles.

\subsection{Universal Computation}

The theoretical foundations of neural computation were established by
\citet{siegelmann1995computational}, who proved that recurrent neural
networks with rational weights are Turing complete. More recently,
\citet{perez2021turing} showed that transformers and Neural GPUs are also
Turing complete, without requiring external memory. These results
establish that neural networks can, in principle, simulate any computable
function, but they do not address the question of how such networks can be
discovered or evolved.

\section{The Present Work}

In this paper, we present \modnet{}, a self-growing layer-free recurrent
network that addresses the limitations of fixed architectures by
discovering its own structure through evolution. The architecture has
four defining characteristics:

\begin{enumerate}
  \item \textbf{Layer-free.} There are no predefined layers, and no
    modular grouping. Every node is a general-purpose threshold unit that
    may connect to any other node, including itself.

  \item \textbf{Outputs as ordinary nodes.} The first $n_{\text{out}}$
    nodes of the network are designated as outputs; there is no separate
    output layer.

  \item \textbf{Grown and pruned.} The network starts from a small random
    topology and grows through addition of nodes and pruned through removal
    of inactive or unreachable ones. Growth is bidirectional.

  \item \textbf{Gradient-free.} No gradient information is used at any
    stage. All learning is performed by a compact evolutionary algorithm
    operating on the topology and weights simultaneously.
\end{enumerate}

We evaluate \modnet{} on a suite of \textbf{sixteen} benchmark problems
spanning four broad classes: Boolean logic, digital arithmetic,
sorting networks, and continuous regression. Table~\ref{tab:intro-headline}
summarizes the headline results.

\begin{table}[h]
  \centering
  \small
  \caption{Headline results of \modnet{} on the sixteen benchmark
    problems. Boolean tasks are measured in exact accuracy; continuous
    tasks in root-mean-square error (RMSE). Full details are reported
    in Chapter~\ref{ch:results}.}
  \label{tab:intro-headline}
  \begin{tabular}{llrrr}
    \toprule
    \textbf{Task} & \textbf{Type} & \textbf{Result} & \textbf{Nodes} & \textbf{Time (s)} \\
    \midrule
    XOR-2               & Boolean    & \textbf{100.0\%} &  8 &  0.8 \\
    XOR-3               & Boolean    & \textbf{100.0\%} &  9 &  2.0 \\
    Parity-4            & Boolean    & \textbf{100.0\%} & 10 &  1.8 \\
    Parity-6            & Boolean    & 95.3\%           &  8 & 40.6 \\
    Majority-3          & Boolean    & \textbf{100.0\%} & 26 &  0.9 \\
    Majority-5          & Boolean    & \textbf{100.0\%} & 26 & 10.2 \\
    Full Adder          & Boolean    & \textbf{100.0\%} & 26 &  3.0 \\
    2-bit Adder         & Boolean    & \textbf{100.0\%} & 26 &  7.3 \\
    Comparator          & Boolean    & \textbf{100.0\%} & 27 &  3.8 \\
    MUX 4-to-1          & Boolean    & \textbf{100.0\%} & 26 &  4.8 \\
    Sort-4              & Boolean    & \textbf{100.0\%} & 26 &  8.7 \\
    \midrule
    $\sin(2\pi x)$      & Regression & RMSE $= 0.0465$  & 18 & 30.0 \\
    $\sin(4\pi x)$      & Regression & RMSE $= 0.0789$  & 13 & 25.1 \\
    $x^2 + y^2$         & Regression & RMSE $= 0.0143$  & 26 & 21.1 \\
    Gaussian            & Regression & RMSE $= 0.0275$  & 27 & 41.6 \\
    Multiplication      & Regression & RMSE $= 0.0254$  & 13 & 28.5 \\
    \bottomrule
  \end{tabular}
\end{table}

Three observations follow immediately from
Table~\ref{tab:intro-headline}. First, \modnet{} solves \textbf{fourteen
of sixteen} tasks to their success threshold, including all of the digital
arithmetic and sorting tasks. Second, the evolved networks are small: no
task required more than $27$ nodes, and the multiplication task required
only $13$. Third, the same architecture, without any task-specific
modification, handles Boolean logic, digital arithmetic, sorting, and
continuous regression.

\section{Key Findings}

Beyond the raw numbers, we observe several structural properties that
emerge spontaneously from the evolutionary process.

\subsection{Recurrence Emerges, and Scales with Complexity}

Every successfully evolved network contains recurrent self-loops, ranging
from one self-loop (XOR-2) to fourteen (Gaussian). No mechanism in our
algorithm explicitly favors recurrence. Self-loops emerge because the
unrestricted topology allows them, and evolutionary pressure discovers them
to be useful.

Crucially, the number of self-loops scales with the temporal complexity of
the task. XOR-2, a purely combinatorial problem, requires only one
self-loop. Gaussian and multiplication, both continuous, require ten to
fourteen. This suggests that recurrence is a \emph{structural necessity}
for tasks that require temporal integration, and that \modnet{} discovers
this necessity without being told.

\subsection{Sparsity Is the Norm}

Evolved networks use a small fraction of the available connections. Density,
defined as the ratio of active connections to total possible connections
($E / N^2$), ranges from $27\%$ to $67\%$ depending on the task. This
sparsity was not imposed by any regularizer; it emerged from the pruning
mechanism, which removes nodes whose activity is either too low or too
high.

The observed sparsity is consistent with that of biological neural
networks \citep{chklovskii2004wiring}, where the vast majority of
potential synapses are absent.

\subsection{Growth Adapts to Task Complexity}

The initial network has $26$ nodes. Over evolution, networks either
shrank, grew, or stayed approximately the same:

\begin{itemize}
  \item \textbf{Shrinking:} XOR-2 ($26 \to 8$), XOR-3 ($26 \to 9$),
    Parity-4 ($26 \to 10$), Multiplication ($26 \to 13$).
  \item \textbf{Growing:} sin(2$\pi x$) ($26 \to 18$), Gaussian
    ($26 \to 27$).
  \item \textbf{Nearly stable:} Parity-6, Majority-5, Full Adder.
\end{itemize}

Boolean tasks tend to converge to small networks, while continuous tasks
require larger networks. The pattern is consistent with the intuition that
Boolean functions have low algorithmic complexity, while continuous
functions require more expressive approximations.

\subsection{Layer-Free and Modular Architectures Are Not Equivalent}

In earlier work we studied a modular variant of \modnet{}, where nodes
were grouped into a fixed number of modules. The present work is fully
layer-free: there are no modules. A direct comparison reveals that the
layer-free variant is \emph{equal or superior} on every benchmark,
including a decisive win on the multiplication task, where it uses only
$13$ nodes versus $18$ in the modular variant.

This suggests that modular scaffolds, far from helping, can \emph{restrict}
the search when the task does not require modular structure.

\section{Contributions}

The main contributions of this work are:

\begin{enumerate}
  \item \textbf{Architecture.} We introduce a fully layer-free, recurrent
    network architecture with no modular grouping and no separate output
    layer. The entire structure, including the number of nodes, is
    discovered through evolution.

  \item \textbf{Comprehensive benchmark suite.} We evaluate on
    \textbf{sixteen} tasks spanning Boolean logic, digital arithmetic,
    sorting networks, and continuous regression. To the best of our
    knowledge, this is the first time a single layer-free evolved
    architecture has been demonstrated on such a range of tasks.

  \item \textbf{Digital circuit results.} We show that \modnet{} solves
    full adder, 2-bit adder, comparator, 4-to-1 multiplexer, and 4-bit
    sorting network, all with exact $100\%$ accuracy. This establishes
    that evolved layer-free networks can implement digital logic of
    nontrivial complexity.

  \item \textbf{Emergent recurrence.} We observe that recurrence emerges
    spontaneously and scales with the temporal complexity of the task.
    XOR-2 requires one self-loop; Gaussian requires fourteen.

  \item \textbf{Comparison with modular variant.} We compare against a
    modular variant of the same algorithm, and find the layer-free
    variant equal or superior on every benchmark.

  \item \textbf{Reproducibility.} We provide a complete, open-source
    implementation and a reproducible benchmark suite, including per-task
    graphs in force-directed (spring) layout, event logs of every birth
    and prune operation, and full training curves.
\end{enumerate}

\section{Organization}

The remainder of this thesis is organized as follows.

\textbf{Chapter~\ref{ch:related}} reviews related work in neuroevolution,
growing networks, modular neural networks, sparse training, and universal
computation.

\textbf{Chapter~\ref{ch:method}} describes the \modnet{} architecture in
detail, including the node model, the layer-free connectivity, the
recurrence mechanism, the fitness function, and the evolutionary
algorithm.

\textbf{Chapter~\ref{ch:experiments}} presents the experimental setup:
the sixteen benchmark problems, the data encoding, the evaluation
procedure, and the hyperparameters.

\textbf{Chapter~\ref{ch:results}} reports the results, including the
summary table, per-task network graphs in spring layout, learning curves,
and a detailed analysis of the two failure cases (Parity-6 and
$\sin(4\pi x)$).

\textbf{Chapter~\ref{ch:discussion}} interprets the findings, discusses
limitations and threats to validity, and proposes directions for future
work.

\textbf{Chapter~\ref{ch:conclusion}} concludes the thesis.

\textbf{Appendix~\ref{app:hyperparameters}} lists all hyperparameters.
\textbf{Appendix~\ref{app:snapshots}} presents network snapshots at
selected generations. \textbf{Appendix~\ref{app:algorithm}} provides the
complete algorithm in pseudocode.

% ============================================================
%  Bibliography for this chapter
% ============================================================
